% Draft rewrite snippets for reviewer-driven revision.
% Target: paper_siga/main.tex abstract, introduction opening/contributions,
% and Method Sections 3.2--3.4.
%
% Integration note: this file is intentionally not included by main.tex.
% It is a staging draft only.
%
% Suggested old paragraphs to comment out and preserve before replacement:
% 1. Current abstract: preserve as commented text. It contains useful experiment
%    claims, but the narrative should move away from "procedural + generative
%    bridge" and remove draft/status phrases such as "technically compatible".
% 2. Introduction paragraphs 1--2: preserve as background material for Related
%    Work or the second half of the introduction. They are useful, but should no
%    longer be the opening motivation.
% 3. Introduction paragraph beginning "Our experiments indicate that naive
%    bridges fail.": preserve the substance, but make it the causal opening
%    of the new story.
% 4. Current Contributions item list: preserve as a checklist. Replace the main
%    text with a stronger language/algorithm/projection/material framing.
% 5. Method 3.2 "Typed Sparse-Latent State": comment out and preserve for an
%    appendix or implementation note. The main text should use the minimal state
%    z_d=(V_d,F_d,A_d), not the eight-component state tuple.
% 6. Method 3.3 "Grammar and Rules": comment out and preserve for appendix
%    detail. Replace the giant grammar tuple with objects of the language and
%    a compact rule template.
% 7. Method 3.4 "Operational Semantics": comment out and preserve the merge
%    equations if desired, but replace the long operator chain with the compact
%    recursive transition plus Algorithm 1 and Algorithm 2.
% 8. Method 3.5 "Classical Systems as Limits": preserve, but move after the core
%    algorithm or to appendix. In the main method, one short paragraph is enough.

\section*{Revised Abstract Draft}

Native 3D generators increasingly expose sparse latent representations that
look well suited for recursive structure: a tree, root system, coral, ornament,
or architectural motif can in principle be grown by repeatedly editing local
3D tokens. In practice, direct uses of such generators are unstable. A 2D
scaffold can be ignored or inflated into sheets, direct sparse edits create
floating fragments, global flow repair can wash out the intended topology, and
final-only cleanup cannot prevent fragments from becoming new growth roots at
the next recursion depth. We introduce a generation-model-native recursive
language over sparse 3D latents. A program state consists of active sparse
tokens, their latent features, and grammar-readable anchors. Rule templates
propose local growth or transform-copy edits on this state, while a frozen
Trellis2 model is used only for masked local naturalization, decoding,
re-encoding, and texture export. The key execution principle is projection
inside the recursive transition: after each depth, the decoded candidate is
projected to a connected, admissible asset and re-encoded before further rules
can fire. This turns connected support into a recursive invariant rather than a
post-hoc mesh cleanup target. Experiments on branching, frontier-growth, and
transform-copy programs show that per-depth projection stabilizes finite-depth
recursive growth, exposes a stability-expression tradeoff across rule families,
and supports textured GLB export for selected projected meshes without training
a new 3D generator.

\section*{Revised Introduction Opening Draft}

Sparse latent 3D generators suggest a tempting route to recursive modeling.
Instead of asking a model to synthesize a complete object in one shot, one could
grow a tree, root, coral, ornament, or architectural asset by repeatedly editing
only the currently active 3D tokens. This would make the generator's native
state space the substrate of a recursive program: topology would be controlled
by rules, while local geometry and appearance would be supplied by a learned
asset prior.

The direct version of this idea fails for a specific reason. Native 3D
generators are strong local and object-level priors, but they are not recursive
execution engines. When a procedural scaffold is rendered as a 2D condition,
image-to-3D generation often produces sheets, blobs, or disconnected fragments
instead of preserving the intended branching topology. When sparse coordinates
are edited directly, small floating components accumulate across depth. When a
global flow sampler is used as a repair step, it can replace the grammar's
topology with the generator's preferred object-level shape. When cleanup is
deferred until the final mesh, invalid fragments may already have acted as
parents, frontiers, cached motifs, or material sources in intermediate states.

These failures point to a different formulation. The generator should not own
the global recursive structure. Instead, the recursive program should own the
sparse support, anchors, and attachment graph, and the frozen generator should
act as a masked local prior around rule-proposed edits. We therefore formulate
recursive asset synthesis as a \emph{generation-model-native recursive language
over sparse 3D latents}. The language executes over the same kind of sparse
state used by the 3D generator, but its semantics are grammar-first: rules
select handles, propose local token edits, preserve attachment certificates,
and project each decoded candidate back to an admissible connected asset before
the next depth.

This framing differs from both classical procedural modeling and one-shot 3D
generation. Classical grammars expose recursion and control, but typically
output curves, tubes, voxels, or hand-designed meshes that lack learned surface
and material priors. One-shot 3D generators produce rich assets, but provide no
native guarantee that a multi-depth recursive program remains attached,
addressable, and editable. Our goal is narrower and more testable than
unbounded world generation: finite-depth recursive assets whose structure is
controlled by a sparse-latent grammar and whose local realization is supported
by a frozen 3D generator.

\paragraph{Contributions.}
This paper makes the following contributions:
\begin{itemize}
  \item We introduce a generation-model-native recursive language over sparse
  3D latents, where programs operate on sparse token support, latent features,
  and grammar-readable anchors rather than on 2D scaffolds or post-hoc meshes.
  \item We define a compact rule-template semantics for recursive growth,
  transform-copy, frontier accretion, and local refinement over sparse latent
  states, avoiding a monolithic system tuple.
  \item We propose projection-stabilized execution: each recursion step performs
  rule proposal, masked local naturalization, decoding, connected projection,
  and re-encoding before later rules can use the state.
  \item We give an explicit projection algorithm for connected admissible assets
  and show why per-depth projection provides a recursive invariant that
  final-only cleanup lacks.
  \item We evaluate finite-depth recursive assets across branching,
  frontier-growth, and transform-copy rule families, and report texture/PBR
  export compatibility for selected projected meshes using a frozen Trellis2
  pipeline.
\end{itemize}

\section*{Method Rewrite Draft: Sections 3.2--3.4}

\subsection*{3.2 Core State and Handles}

We write the recursive state at depth \(d\) as the minimal sparse-latent object
\begin{equation}
z_d=(\mathcal V_d,\mathbf F_d,\mathcal A_d).
\end{equation}
Here \(\mathcal V_d\subset h\mathbb Z^3\) is the active sparse token support,
\(\mathbf F_d:\mathcal V_d\rightarrow\mathbb R^q\) stores the generator-native
latent features attached to those tokens, and \(\mathcal A_d\) is the
grammar-readable anchor structure. The anchors include typed handles, local
frames, ownership regions, attachment edges, frontier labels, and optional
material handles. A handle can be written as
\begin{equation}
h_i=(\sigma_i,T_i,\Omega_i,\alpha_i),
\end{equation}
where \(\sigma_i\) is its type, \(T_i\) is a local frame, \(\Omega_i\subseteq
\mathcal V_d\) is its owned support or target region, and \(\alpha_i\) stores
local attributes such as radius, age, material intent, or rule priority.

This minimal state separates the semantics from implementation bookkeeping.
Masks, random seeds, rule traces, material export metadata, diagnostic scores,
and cached motifs are auxiliary data used by the executor, but they are not the
primary mathematical state a reader must remember. The important invariant is
that each active handle must be addressable in the sparse support and attached
to the recursive scaffold unless it is explicitly marked as an inactive
descriptor. Thus a small floating component is not merely a visual artifact: if
it enters \(\mathcal A_d\) as an active handle, it can contaminate all later
recursion depths.

For this reason we maintain a connected-state domain. Let \(R_d\subset
\mathcal V_d\) denote root anchors or scaffold tokens, and let
\(G_\eta(\mathcal V_d)\) be the \(\eta\)-neighborhood graph over sparse tokens.
The attached support is
\begin{equation}
\mathcal V_{\mathrm{att}}(z_d)=
\{v\in\mathcal V_d:\ v\leadsto R_d\ \mathrm{in}\ G_\eta(\mathcal V_d)\}.
\end{equation}
We use the violation
\begin{equation}
V_{\mathrm{conn}}(z_d)=
1-\frac{\operatorname{mass}(\mathcal V_{\mathrm{att}}(z_d))}
{\operatorname{mass}(\mathcal V_d)+\epsilon}.
\end{equation}
An admissible recursive state satisfies bounded connectivity, occupancy,
renderability, material-seam, and token-budget violations:
\begin{equation}
\mathcal Z_{\mathrm{adm}}(\lambda)=
\{z:\ V_{\mathrm{conn}}(z)\le\epsilon_c,\ 
V_{\mathrm{frontier}}(z)=0,\ 
V_{\mathrm{occ}}(z)\le\epsilon_o,\ 
V_{\mathrm{mat}}(z)\le\epsilon_m,\ 
|\mathcal V(z)|\le T_{\max}\}.
\end{equation}
The condition \(V_{\mathrm{frontier}}(z)=0\) means that every active frontier
or growth handle has a valid attachment path to the root scaffold. Inactive
LOD descriptors or cached motifs may be stored by the implementation, but they
cannot generate new geometry until projection certifies their attachment.

\subsection*{3.3 Rule Templates as a Recursive Language}

The language consists of three objects: typed handles, rule templates, and
realization operators. A rule is not a giant grammar tuple; it is a local
rewrite template over handles:
\begin{equation}
\rho:\ h_i\mapsto
\{(\sigma_j,T_j,\Omega_j,m_j,\kappa_j,\rho_j)\}_{j=1}^{k}.
\end{equation}
Each emitted item creates or updates a handle of type \(\sigma_j\), with local
transform \(T_j\), target region \(\Omega_j\), edit mask \(m_j\), attachment or
feasibility condition \(\kappa_j\), and proposal kernel \(\rho_j\). The
proposal kernel may instantiate a branch segment, extend a frontier, copy a
motif through a transform, split a local region, densify a terminal detail, or
request a masked generative naturalization step.

Rules are executed only when their feasibility conditions hold. In particular,
a materialized proposal \(\Delta\mathcal V\) must lie in the projectable
reachable domain or carry a bridge certificate:
\begin{equation}
\Delta\mathcal V\subset
\mathcal R_{\mathrm{proj}}(z_d;b_d,r_d)
\quad\mathrm{or}\quad
\exists B:\ \Delta\mathcal V\cup B\in\mathcal Z_{\mathrm{adm}}(\lambda_d),
\ \operatorname{cost}(B)\le b_d.
\end{equation}
Here \(\mathcal R_{\mathrm{proj}}\) contains sparse coordinates that are close
enough, collision-free enough, and cheap enough to connect back to the attached
scaffold. This condition is the semantic difference between recursive growth
and unconstrained latent editing: a patch is not active merely because it looks
plausible locally; it must remain usable by the next recursive step.

\begin{table}[t]
  \centering
  \caption{Representative rule families in the sparse-latent recursive
  language. Each family uses the same handle-to-proposal template but different
  proposal kernels and feasibility conditions.}
  \label{tab:rule-families-draft}
  \begin{tabular}{lll}
    \toprule
    Rule family & Proposal kernel & Typical use \\
    \midrule
    grow & extend attached frontier & vines, roots, branches \\
    branch & split one handle into several attached handles & trees, bushes \\
    attach & accept frontier hit with bridge certificate & porous and DLA-like growth \\
    copy & transform a connected motif & ornaments, portals, architecture \\
    refine & densify a local terminal region & high-detail recursive tips \\
    material & transport or blend material handles & texture/PBR continuity \\
    \bottomrule
  \end{tabular}
\end{table}

Given a selected rule set \(\mathcal R_d\), rule proposal and sparse merge are
summarized as
\begin{equation}
\tilde z_{d+1}=\mathcal M(z_d,\operatorname{Prop}_{\mathcal R_d}(z_d)).
\end{equation}
For a transform-copy proposal, the support and features are
\begin{equation}
\widehat{\mathcal V}_j=Q_h(T_j(\mathcal V_d\cap\Omega_j)),\qquad
\widehat{\mathbf F}_j(Q_h(T_jv))=\mathcal T_j^F\mathbf F_d(v),
\end{equation}
and masked feature merge is
\begin{equation}
\mathbf F^+(v)=
\frac{w_0(v)\mathbf F_d(v)+\sum_j w_j(v)\widehat{\mathbf F}_j(v)}
{w_0(v)+\sum_jw_j(v)+\epsilon},
\qquad
w_j(v)=m_j(v)\beta_j(v).
\end{equation}
Material handles use the same transport-and-merge logic, but are kept distinct
from the geometric claim: current evidence supports texture/PBR export for
selected projected meshes, while per-depth material naturalization remains a
controlled extension of the same semantics.

\subsection*{3.4 Projection-Stabilized Execution}

The recursive transition is intentionally compact:
\begin{equation}
z_{d+1}=
\operatorname{Enc}_{\theta}
\left[
\Pi_{\lambda_d}
\left(
\operatorname{Dec}_{\theta}
\left(
\mathcal T_{\theta}(z_d;\mathcal R_d,y)
\right)
\right)
\right].
\label{eq:compact-recursive-transition-draft}
\end{equation}
The operator \(\mathcal T_{\theta}\) contains rule selection, sparse proposal,
masked merge, feasibility filtering, and optional local naturalization. The
decoder \(\operatorname{Dec}_{\theta}\) maps the sparse state to a mesh or
occupancy proxy, \(\Pi_{\lambda_d}\) projects the decoded candidate to an
admissible connected asset, and \(\operatorname{Enc}_{\theta}\) re-encodes the
projected asset before the next depth. Projection is therefore inside the
recursive map. It is not a final cleanup pass.

\begin{algorithm}[t]
\caption{Recursive Sparse-Latent Program Execution}
\label{alg:rslg-execution-draft}
\begin{algorithmic}[1]
\REQUIRE root mesh \(x_0\), recursive program \(\mathcal G\), frozen generator
\(\theta\), condition \(y\), depth \(D\), projection schedule
\(\{\lambda_d\}_{d=0}^{D-1}\)
\STATE \(z_0\leftarrow \operatorname{Enc}_\theta(x_0)\)
\STATE initialize anchors \(\mathcal A_0\), root handles, material handles, and
auxiliary caches
\FOR{\(d=0,\ldots,D-1\)}
  \STATE select active handles \(H_d\subset\mathcal A_d\)
  \STATE select applicable rules \(\mathcal R_d\leftarrow
  \operatorname{SelectRules}(H_d,\mathcal G,z_d)\)
  \STATE generate sparse proposals
  \(P_d\leftarrow \operatorname{Prop}_{\mathcal R_d}(z_d)\)
  \STATE reject or defer proposals outside \(\mathcal R_{\mathrm{proj}}\) unless
  they carry a valid bridge certificate
  \STATE merge accepted proposals with masks and feature/material blend kernels:
  \(\tilde z_{d+1}\leftarrow\mathcal M(z_d,P_d)\)
  \STATE naturalize edited regions with the frozen masked flow:
  \(\bar z_{d+1}\leftarrow\mathcal N_\theta(\tilde z_{d+1};m_d,y)\)
  \STATE decode candidate \(x_{d+1}\leftarrow
  \operatorname{Dec}_\theta(\bar z_{d+1})\)
  \STATE project connected admissible asset
  \(x_{d+1}^{\star}\leftarrow
  \Pi_{\lambda_d}(x_{d+1},\mathcal A_d)\)
  \STATE re-encode \(z_{d+1}\leftarrow
  \operatorname{Enc}_\theta(x_{d+1}^{\star})\)
  \STATE update anchors, frontiers, material handles, and connected caches
\ENDFOR
\RETURN \(x_D^\star=\operatorname{Dec}_\theta(z_D)\) or exported textured asset
\end{algorithmic}
\end{algorithm}

\paragraph{Masked flow naturalization.}
The frozen generator is used as a local prior under a hard grammar mask. In
deterministic flow form,
\begin{equation}
\frac{d z_t}{dt}=v_\theta(z_t,t,y).
\end{equation}
Let \(m_d\) be the sparse edit mask selected by the rule and clipped by
connectivity feasibility. A single reverse step is clamped as
\begin{equation}
z_{t-\Delta t}=
(1-m_d)\odot z_{\mathrm{anchor}}
+m_d\odot\Phi_{\theta,t}(z_t,y).
\label{eq:masked-flow-clamp-draft}
\end{equation}
Equivalently, for sparse features on edited tokens,
\begin{equation}
\mathbf f_i'=
(1-\alpha_i)\mathbf f_i^{\mathrm{rule}}
+\alpha_i\mathbf f_i^{\theta},
\qquad i\in m_d.
\end{equation}
The anchor region is hard-clamped, and \(m_d\) is zero outside the grammar's
target region and outside the projectable reachable domain. The sampler
therefore naturalizes local shape or material; it does not decide the global
recursive topology.

\begin{algorithm}[t]
\caption{Projection to a Connected Admissible Asset}
\label{alg:projection-draft}
\begin{algorithmic}[1]
\REQUIRE decoded candidate \(x\), previous anchors \(\mathcal A_d\), root set
\(R_d\), projection parameters \(\lambda_d\)
\STATE build a voxel, occupancy, or mesh connectivity graph \(G(x)\)
\STATE identify root-attached components by graph reachability from \(R_d\)
\STATE score each component by mass, attachment distance, bridge cost,
renderability, material consistency, and rule ownership
\STATE keep all valid root-attached components and any component with an
accepted bridge path under budget \(b_d\)
\STATE prune orphan components that cannot be bridged without violating
\(\lambda_d\)
\STATE optionally weld, close small gaps, remesh, or insert bridge geometry for
accepted components
\STATE update handles: remove orphan frontiers, reattach surviving handles to
projected support, and mark uncertain motifs inactive
\STATE validate token budget, connected-frontier invariant, material seams, and
renderability
\RETURN projected asset \(x^\star\) and updated anchors \(\mathcal A^\star\)
\end{algorithmic}
\end{algorithm}

Algorithm~\ref{alg:projection-draft} is written at the asset/proxy level because
the available implementation may use voxel connectivity, mesh components, or an
occupancy proxy before re-encoding. A bridge-aware version generalizes
prune-only projection. For a candidate component \(Q\), define bridge paths
\begin{equation}
\Gamma(Q,R_d)=
\{\gamma=(q_0,\ldots,q_T):q_0\in Q,\ q_T\in R_d,\ 
q_{t+1}-q_t\in\mathcal N_\eta\}.
\end{equation}
The projection objective can be expressed as
\begin{equation}
\begin{split}
\Pi_{\lambda}(x)=
\operatorname*{arg\,min}_{y\in\mathcal X_{\mathrm{adm}}(\lambda)}
&\ d_{\mathcal X}(x,y)
+\eta_{\mathrm{obs}}d_{\mathrm{obs}}(x,y)
+\eta_m d_{\mathrm{mat}}(x,y)\\
&+\lambda_{\mathrm{del}}\operatorname{mass}(R)
+\lambda_{\mathrm{br}}\sum_{\gamma\subset B}\operatorname{len}(\gamma)
+\lambda_{\mathrm{front}}V_{\mathrm{frontier}}(y),
\end{split}
\end{equation}
where \(R\) is removed orphan support and \(B\) is inserted or certified bridge
support. Prune-only projection is the limiting case
\(\lambda_{\mathrm{br}}\rightarrow\infty\). This distinction matters for DLA,
crystal, radial, and fork-like rules: deleting every thin or distant component
can destroy legitimate recursive structure, while accepting every component
creates active floating frontiers. Projection resolves the conflict by making
attachment a per-depth admissibility test.

\paragraph{Why final-only cleanup is insufficient.}
Let \(e(z)\) combine connectivity, orphan-frontier, renderability, and material
violations. If each projection satisfies
\begin{equation}
e(\Pi_{\lambda_d}(x))\le\epsilon_P
\end{equation}
and the decode/re-encode round trip adds at most \(\epsilon_E\), then the
executed recursive states obey
\begin{equation}
e(z_d)\le \epsilon_P+\epsilon_E,\qquad d=1,\ldots,D.
\end{equation}
Final-only cleanup lacks this invariant because an invalid component can be
selected as a parent, copied as a motif, sampled as a local condition, or stored
in a cache before the final mesh is cleaned. This is the central algorithmic
reason to place projection inside Eq.~\eqref{eq:compact-recursive-transition-draft}.

\paragraph{Placement of classical systems.}
Classical procedural systems can be mentioned after the algorithm as restricted
rule families: L-systems use ordered handles and synchronous branch rules,
space colonization uses attractor-guided grow proposals, DLA uses frontier
hitting proposals with attachment certificates, and IFS-like assets use
transform-copy rules. The main text should not expand these into a separate
large formal detour unless space allows; detailed embeddings are better suited
for the appendix.
